%=====================================================================
\section{Related work: what the papers did}\label{sec:literature}
%=====================================================================
This section describes the papers our work builds on, what each one contains, and how we use it. We searched web search engines, Semantic Scholar, OpenAlex, Crossref and arXiv, and followed the citation trails of the key papers (about 620 citing titles screened). We say how deeply we read each paper, because a claim can only be as reliable as the reading behind it. ``Full'' means the whole text was read; ``partial'' means the relevant sections; ``abstract'' means only the abstract or descriptions by citing papers.

\subsection{The three papers we validate against}

\paragraph{Mueller and D'Andrea, ICRA 2014~\cite{mueller2014stability} (read in full).}
This is the paper our project rests on. It contains:
\begin{itemize}
  \item \textbf{A model} of a quadrotor (their eqs.~1 to 14) that includes the gyroscopic effect of the spinning propellers (explained in Section~\ref{sec:model}) and a linear yaw drag torque $-\gamma r$. The propeller's own angular acceleration term is neglected.
  \item \textbf{Vehicle parameters} (their Sec.~V-A; the motor time constant is in Sec.~V-C), based on an Ascending Technologies Hummingbird (a commercial research quadrotor): mass 0.50\,kg, arm length 0.17\,m, inertias $I^T_{xx}=3.2\times10^{-3}$ and $I^T_{zz}=5.5\times10^{-3}$\,kg\,m$^2$, propeller inertia $1.5\times10^{-5}$\,kg\,m$^2$ (printed as ``kg\,mm$^2$'' in the paper, an evident typo), thrust coefficient $6.41\times10^{-6}$\,N\,s$^2$/rad$^2$, torque-to-thrust ratio 0.0169\,m, yaw drag $2.75\times10^{-3}$\,N\,m\,s/rad, thrust range 0.2 to 3.8\,N per rotor, and a motor time constant of 15\,ms. We use every one of these values.
  \item \textbf{The relaxed-hover equilibria} for the loss of one propeller (their eqs.~50 to 53: thrusts (2.05, 1.02, 2.05, 0)\,N, body rates (0, 5.69, 18.89)\,rad/s, axis (0, 0.289, 0.958), orbit radius 8\,mm for the thrust ratio 0.5), of two opposing propellers (eqs.~57 to 59) and of three propellers (eqs.~62 to 64), plus a figure (Fig.~3) of how the one-rotor equilibrium changes with the thrust ratio.
  \item \textbf{A linear model} of the reduced attitude about the equilibrium, the statement that it is controllable (for one lost propeller) when $l\ne0$ and $\bar n_z\ne0$, and the coupling constant $\bar a$ (their eq.~28) that decides when controllability is lost (eqs.~37 to 39).
  \item \textbf{An LQR controller} for one lost propeller with weights 1 on the rates, 20 on the reduced attitude, 0 on the motor states and 1\,N$^{-2}$ on the inputs, with motor dynamics included as extra states (the two-propeller design uses different weights).
  \item \textbf{A flight experiment} (their Fig.~4): propeller 4 is switched off at $t=0$ \emph{with the failure known}. The drone first spins up on the two opposing propellers 1 and 3, and the LQR is switched on once the rotation rate exceeds 10\,rad/s, which took 0.62\,s. Reading the plot, during spin-up rotors 1 and 3 run at a constant thrust of about 2.95\,N with rotor 2 stopped; the drone then climbs from 2.0 to about 2.9\,m before settling a few tenths of a metre from its set-point (the paper does not discuss this offset). The measured spin in hover was $\vect\omega = (2.1, 3.7, 13.7)$\,rad/s (their eq.~55), about 14.3\,rad/s against about 19.7\,rad/s predicted. In a separate two-propeller flight (their Fig.~5) the drone stayed about 2.4\,m below its set-point, which the authors attribute to the propellers producing ``significantly less thrust than expected when the body is rotating at high angular velocities''.
  \item \textbf{Future work} stated by the authors: ``implement the presented results with a fault detection scheme, and investigate different switching strategies for transitioning from e.g.\ four propellers to three''. This is exactly where RQ2 sits.
\end{itemize}
\emph{Our use:} the model, all parameters, the equilibria and Fig.~3 as validation targets, the LQR weights, the spin-up strategy and the flight spin-up time as a flight-data check; the thrust loss seen in the two-propeller flight motivates our truth plants (Section~\ref{sec:truth}).

\paragraph{Mueller and D'Andrea, IJRR 2016~\cite{mueller2016relaxed} (Sections 1 to 5 read).}
A more general and more detailed version for any multicopter. Two physical effects are added: (i) the thrust of a propeller depends on its speed relative to the \emph{air}, $\Omega_i + r$, not relative to the body, because the whole body spins at rate $r$; (ii) the drag torque is \emph{quadratic} in the body rate, $-\|\vect\omega\|K_d\vect\omega$, on all three axes. The paper expresses its hover solutions in propeller speeds and lists, in its Sec.~5.2 (eqs.~31 to 50), the hover solutions (propeller speeds and body rates) for five cases: no failure, one rotor lost, two opposing rotors lost, two adjacent rotors lost and three rotors lost. In flight with one rotor lost, the measured spin rate was 15.8\,rad/s against 19.9\,rad/s in simulation (Sec.~5.4); the authors suggest additional torques from the translation of the propeller centres as a possible cause.
\emph{Our use:} we wrote an independent implementation of this model and checked it against the five published solutions.

\paragraph{Sun, Wang, Chu and de Visser, T-RO 2021~\cite{sun2020incremental} (abstract and relevant equations read).}
An incremental nonlinear dynamic inversion (INDI) controller for a Parrot Bebop~2 that loses two opposing rotors, tested in a wind tunnel at speeds above 8\,m/s. An LQR of the Mueller--D'Andrea type is compared and degrades in these conditions. The paper gives a closed-form equilibrium yaw rate for two lost rotors (its eq.~60) and the Bebop~2 parameters (its Table~I in the arXiv version, which the text refers to as Table~III).
\emph{Our use:} a second, different airframe to test our equilibrium solver.

\subsection{Other controllers for complete rotor loss}
\begin{itemize}
  \item \textbf{Freddi, Lanzon and Longhi}~\cite{freddi2011feedback,lanzon2014flight} (abstract): give up yaw and stabilise a constant-rate spin about the vertical axis with level attitude, using feedback linearization and $H_\infty$ loop shaping (a robust linear design method) around this working point.
  \item \textbf{Stephan, Schmitt and Fichter}~\cite{stephan2018lpv} (abstract only): an LPV (linear parameter-varying) controller in which the yaw rate is a time-varying scheduling parameter, with guaranteed performance during the transient after the actuator loss. We have not been able to read the full text, so we do not know whether thrust saturation is treated. This is the closest work to our planned certificate over the spin-up.
  \item \textbf{Khattab, Alwi and Edwards}~\cite{khattab2019lpv} (abstract): an LPV reduced-attitude model with sliding-mode control.
  \item \textbf{Sun, Baert, Strack van Schijndel and de Visser}~\cite{sun2020upset} (partial): recovery from large upsets using a quadratic-programming control allocation and the ``attainable moment set''. Tested by Monte Carlo and by throwing the drone by hand. It states that a drone ``may be upside down and fast rotating before the FTC is triggered, because of \ldots\ delay of the fault detection module'', and that linear controllers ``could fail'' through linearization errors and saturation. This sentence directly motivates RQ2, but the paper does not quantify the delay.
  \item \textbf{Nan, Sun, Foehn and Scaramuzza}~\cite{nan2022nmpc} (partial): nonlinear MPC with per-rotor thrust limits that recovers even from a failure at the top of a flip, in simulation and flight. The authors state that fault detection is outside the scope of the paper; the failed rotor is reported instantly by a mock detector, so there is no fault-detection-delay analysis.
  \item \textbf{Sun, Zhang and Lang}~\cite{sun2022ftmpc} (abstract): nonlinear MPC that learns an unknown complete rotor failure online.
  \item \textbf{Ke, Cai and Quan}~\cite{ke2023uniform} (abstract): a \emph{passive} controller that tolerates one to three rotor failures without any detection step and without switching.
\end{itemize}

\subsection{Guarantees and detection delay}
\begin{itemize}
  \item \textbf{Simha, Vadgama and Raha}~\cite{simha2017geometric} (full): a geometric controller with an analytic region of attraction (a set of starting states that is proven to recover). The fault is assumed to be known from $t=0$; positive thrust is only checked in simulation; no motor lag.
  \item \textbf{Al Lawati and Akhtar}~\cite{allawati2026quasi} (partial, assumptions checked in the text): quasi-static feedback with path invariance and local exponential convergence. They write that they ``assume the availability of a reliable FDI module and an instantaneous transition''.
  \item \textbf{Ba\c{s}ak and Prempain}~\cite{basak2015switching} and \textbf{Ba\c{s}ak's thesis}~\cite{basak2018thesis} (conference abstract and thesis Ch.~6 read): the closest prior work to RQ2. The thesis treats \emph{intermittent} faults, in which a motor (or the control of an actuator) is lost and regained, as a switched linear system (a system that jumps between several linear models). The undetected interval is modelled as ``faulty plant with the healthy controller'' for a fixed time, $T_D = 100$\,ms for the motor-loss case and 500\,ms for loss of control; LMI conditions with several Lyapunov functions guarantee stability as long as switching is not too frequent (an average ``dwell time'' of about 1.1\,s). The delay is fixed in advance rather than maximized, thrust saturation is not included, and the linearizations use Euler angles and the yaw rate instead of the relaxed-hover model.
  \item \textbf{Qi, Theilliol, He and Han}~\cite{qi2016recoverable} (abstract): a ``recoverable set'' (an invariant set of states that can be driven to the equilibrium) at the relaxed hover, with rotor-speed limits, computed by sum-of-squares programming (an extension of LMI methods to polynomial functions). It also covers partial loss of effectiveness and the loss of more than one propeller. Whether it covers the spin-up is unknown to us.
  \item \textbf{Measured tolerances} are empirical: 180\,ms for a variable-pitch ``heli-quad'', found by simulation sweeps~\cite{kulkarni2020heliquad} (partial read; a different vehicle that keeps yaw control). Other papers report the detection delays their controllers worked with, not a maximum: in simulation, nonlinear MPC and INDI recovered with a detection delay of about 0.1\,s~\cite{chen2025learning}, and a nonlinear MPC recovered in hover with 0.21 to 0.22\,s, while success fell to 60 to 90\% at 0.18 to 0.25\,s during trajectory tracking~\cite{zhou2025rotorfailureaware} (full texts read).
  \item \textbf{Detector speed}: Strack van Schijndel, Sun and de Visser~\cite{strack2021fast} report detection of propeller loss in 30 to 130\,ms (26 propeller ejections on a Bebop~2, 95\% interval 28 to 132\,ms); Yang et al.~\cite{yang2026mml1} report about 45\,ms, and Zhou et al.~\cite{zhou2025rotorfailureaware} detect a stopped motor from its measured speed in about 20 to 30\,ms. So published detectors need about 20 to 130\,ms, depending on the method and the sensors.
  \item \textbf{General theory}: set-theoretic fault-tolerant control uses invariant and reachable sets for detection and reconfiguration~\cite{stoican2013set}; Shin and Belcastro~\cite{shin2006performance} compute, with LMIs, a bound on the performance of a fault-tolerant system as a function of the detection time; Rotondo et al.~\cite{rotondo2015virtual,rotondo2016virtual} design controllers that maximize the tolerated delay between a fault and its isolation, with actuator saturation, for linear and LPV systems whose remaining actuators can still compensate the fault; a motion planner based on backward reachability exists for the failed quadrotor after the fault is handled~\cite{ast2024motionplanning}.
\end{itemize}

\subsection{Koopman models and residual learning}
\begin{itemize}
  \item \textbf{EDMD and EDMDc} (partial): Williams, Kevrekidis and Rowley~\cite{williams2015edmd} introduced EDMD, which fits a linear operator on a dictionary of observables by least squares. Korda and Mezi\'c~\cite{korda2018linear} extended EDMD to systems with inputs and showed how the resulting linear predictor is used in MPC; we call this EDMDc (``with control''). With inputs, the lifted linear model is an approximation even for a rich dictionary. Brunton et al.~\cite{brunton2016koopman} discuss which observables make the dynamics linear.
  \item \textbf{Koopman MPC for healthy quadrotors} (abstracts): SE(3) Koopman MPC~\cite{narayanan2023se3} and analytical Koopman lifting~\cite{martini2025kmpc}.
  \item \textbf{Koopman with faults} (abstracts): Lee et al.~\cite{lee2024koopmanftmpc,lee2024koopmanfdi} address loss-of-effectiveness actuator faults (and, for detection, bias faults), not complete rotor loss, with Koopman MPC and Koopman fault detection. The MPC paper updates the input matrix online with a weighted-window EDMD.
  \item \textbf{Physics plus learned Koopman residual}: soft robot arms~\cite{bruder2025residual} (abstract), ground vehicles~\cite{fu2025residual} (abstract), quadrupeds~\cite{narayanan2026rkmpc} (full text; it compares a nominal model, full-state Koopman models and a residual Koopman model, as we do, and reports that a full-state SE(3) Koopman MPC baseline becomes unstable within about 5\,s in simulation), and healthy quadrotors with a deep Koopman lifting that is updated online~\cite{jin2025dekc} (full text). Bruder et al.\ and Fu et al.\ also report that residual models need less training data than purely data-driven ones.
  \item \textbf{Todde, Martini, Rizzo and Valavanis, ICUAS 2026}~\cite{todde2026residual} (read in full): the closest Koopman work. A small neural network (8 inputs, layers of 64 and 32, 4 outputs, 2,788 parameters) predicts the one-step error of an analytical Koopman model of a \emph{healthy} quadrotor. The prediction is evaluated \emph{outside} the MPC optimization, injected as a measured disturbance (so the QP stays convex) and clipped by $d_{\max}$ taken from percentiles of the training errors. Training data: 20 runs of 80\,s with mass, inertia, arm length and motor time constant randomized by $\pm$20\%. In a single simulated stress test with $+30$\% mass, inertia and arm length and $-10$\% thrust and torque coefficients (outside the training range), the overall position RMSE drops from 0.146\,m (Koopman MPC) to 0.095\,m (residual Koopman MPC); under nominal conditions the residual version is worse (0.040 against 0.017\,m). Simulation only, no rotor failure, no guarantees; formal robustness guarantees are listed as future work. \emph{Our use:} we borrow three ideas: residual outside the QP, clipping, and randomized-parameter training data.
  \item \textbf{Koopman with guarantees} (abstract): SafEDMD~\cite{strasser2024safedmd} designs controllers by semidefinite programming with bounds on the EDMD error, which is the starting point for a future certificate of the residual Koopman loop.
\end{itemize}

\subsection{Tools}
Hu and Lin~\cite{hu2001saturation} give LMI conditions for invariant ellipsoids of linear systems with saturating inputs, which we use in Section~\ref{sec:certificate}. Boyd et al.~\cite{boyd1994lmi} is the standard reference for LMIs in control. Practical attitude estimators for the spinning drone use inertial sensors together with vision~\cite{sun2021vision} or a magnetometer~\cite{solanki2022attitude}.

\subsection{The gap}
\begin{enumerate}
  \item We did not find residual Koopman models applied to the complete-rotor-loss, spinning-equilibrium regime (RQ1). The \emph{pattern} of physics plus learned residual exists for healthy drones~\cite{todde2026residual,jin2025dekc}, and a learned failure model inside MPC exists for complete loss~\cite{sun2022ftmpc}; the combination with the relaxed-hover model is what is new.
  \item Several guaranteed recovery methods start their analysis at the moment the failure is known~\cite{simha2017geometric,allawati2026quasi}; passive schemes avoid detection altogether~\cite{ke2023uniform}; Ba\c{s}ak handles a fixed delay without saturation~\cite{basak2018thesis}; and maximum tolerated delays under saturation have been certified for linear systems whose actuators can still compensate the fault~\cite{rotondo2015virtual,rotondo2016virtual}. For complete rotor loss, where the steady state after the failure is a spin, we did not find a certified \emph{maximum} tolerable delay that accounts for thrust saturation and the drift of the undetected phase (RQ2).
\end{enumerate}
We state both as the outcome of our search, not as proof that no such work exists.
